% TOP_PROBLEM: {"problemId": "problem.top500-omission-w045049-absolute-hodge-conjecture", "canonicalTitle": "Deligne's Absolute Hodge Conjecture", "releaseRank": 94, "primaryDomain": "number_theory_arithmetic_geometry", "primaryDomainLabel": "Number theory & arithmetic geometry", "displayStatus": "open_with_solved_subcases", "releaseStatus": "open_with_solved_subcases"}
% !TeX program = lualatex
% Definition and sources only; this file is not an LLM solution attempt.
\documentclass[11pt]{article}
\usepackage[margin=25mm]{geometry}
\usepackage{fontspec}
\setmainfont{Latin Modern Roman}
\usepackage{amsmath,amssymb,mathtools}
\usepackage{unicode-math}
\setmathfont{Latin Modern Math}
\usepackage{xurl,hyperref}
\hypersetup{colorlinks=true,urlcolor=blue}
\setcounter{secnumdepth}{0}
\setlength{\parindent}{0pt}
\setlength{\parskip}{5pt}
\setlength{\emergencystretch}{3em}
\begin{document}
\section*{\texorpdfstring{Deligne's Absolute Hodge Conjecture}{Deligne's Absolute Hodge Conjecture}}\label{94}
\subsection{Definitions and mathematical statement}
A rational Hodge class $\alpha\in H^{2p}(X,{\mathbb{Q}}(p))$ becomes type $(0,0)$ in the twisted Hodge structure. Via algebraic de Rham comparison it has a conjugate on $X^\sigma$ for each $\sigma\in\operatorname{Aut}({\mathbb{C}})$. Absoluteness means
\[
 \forall\sigma\in\operatorname{Aut}({\mathbb{C}}),\qquad
 \alpha^\sigma\in H^{2p}(X^\sigma,{\mathbb{Q}}(p))\text{ is a Hodge class}.
\]
The conjecture requires this for every rational Hodge class on every smooth projective complex variety. Arbitrary field automorphisms of ${\mathbb{C}}$, not only continuous ones or complex conjugation, are quantified.
\par\smallskip\textit{Formulation and concept references:} \href{https://www.math.ens.psl.eu/~charles/AH.pdf}{[S1]}.

\subsection{Short English statement}
Does every rational Hodge class remain rational and of Hodge type after every automorphism of the complex numbers?
\subsection{Sources}
\begin{itemize}
\item[S1]F. Charles and C. Schnell, Notes on absolute Hodge classes (2011), Section 2.5, Conjecture 25.
\url{https://www.math.ens.psl.eu/~charles/AH.pdf}.
\textit{Formulation source.}
\end{itemize}
\end{document}
